% ---------------------------------------------------------------------
%  RESUMO (NBR 6028: parágrafo único, 150 a 500 palavras, voz ativa,
%  terceira pessoa). Reescrito em nov. 2026 com os resultados.
% ---------------------------------------------------------------------
\setlength{\absparsep}{18pt}
\begin{resumo}
Os quadrirrotores são plantas subatuadas, não lineares, acopladas e
instáveis em malha aberta, e o seu voo depende de uma malha de controle de
atitude executada em tempo real. Nos pilotos automáticos de código aberto,
essa malha é feita por controladores proporcional-integral-derivativo (PID)
em cascata, sintonizados de forma empírica ou automática, e não a partir de
um modelo da aeronave. Este trabalho projeta três controladores de atitude
baseados em modelo --- LQI, LQR com observador de estado estendido e LQG
---, implementa-os no \textit{firmware} ArduPilot pelo \textit{framework}
\codigo{AC\_CustomControl} e os compara ao PID nativo, sintonizado pelo
AutoTune. A comparação é feita em simulação \textit{software-in-the-loop}
(SITL), em que o \textit{firmware} compilado voa contra um modelo de dinâmica
de voo. O modelo da aeronave vem das equações de Newton-Euler, é linearizado
em torno do voo pairado e é confrontado com o simulador, no tempo e na
frequência. O LQI e o LQR com observador usam a banda do nativo; o LQG usa um
filtro de Kalman que modela o caminho de medida do \textit{firmware} e é
projetado para não passar aos motores mais ruído do giroscópio que o nativo.
Os ensaios aplicam a todos os controladores, nas mesmas condições, degraus
de atitude e perturbações de torque em impulso, em degrau curto, em degrau
sustentado, em varredura e estocásticas. Há também cenários com ruído de
giroscópio, erro de modelo, falha de motor, saturação e ângulos grandes,
num lote de 558 voos com o \textit{firmware} congelado. Em rolagem e
arfagem, o LQI e o LQR com observador não superam o nativo no pico da
perturbação, e o LQI equivale à cascata nativa com outros ganhos. O LQG é o
único que reduz o pico, de 7\,\% a 11\,\%, também com ruído de giroscópio e
em voos repetidos. A ablação mostra que essa vantagem vem de modelar o
filtro do giroscópio e o erro do ângulo estimado. Em guinada, todos os
propostos eliminam a recuperação lenta do nativo, de \SI{22}{\second} para
0,7 a \SI{2.3}{\second}, mas o próprio PID ressintonizado também o faz; o
pico só cai com mais banda e mais ruído no esforço ou, no LQG, em 7\,\% a
8\,\%. A vantagem no rastreamento do degrau de atitude vem da alimentação
direta pela inversão do modelo, que o próprio PID poderia usar. Os limites
encontrados vêm mais da informação disponível --- o filtro e o ruído do
giroscópio e a baixa autoridade em guinada --- que da lei de controle.

\textbf{Palavras-chave}: quadrirrotor; controle de atitude; ArduPilot;
regulador linear quadrático; filtro de Kalman; simulação
\textit{software-in-the-loop}.
\end{resumo}
